\documentclass[10pt,a4paper]{amsart}
\usepackage[margin=19mm]{geometry}
\usepackage{amsmath,amssymb,mathtools}
\usepackage[scr=rsfs,cal=euler]{mathalfa}
\usepackage{xfrac}
\usepackage[unicode,hidelinks,bookmarks=false]{hyperref}
\newcommand{\ie}{i.e.\ }
\newcommand{\cf}{cf.\ }
\newcommand{\resp}{respectively\ }
\newcommand{\Subsection}{\subsection}
\providecommand{\nobreakdash}{\nobreak\hbox{-}}
\setlength{\emergencystretch}{2em}
\multlinegap0pt
\usepackage{amssymb}
\usepackage[shortlabels]{enumitem}
\usepackage{tikz}
\usepackage{mathtools}
\usepackage{xcolor}
\newtheorem{theorem}{Theorem}
\def\Ext{\mathrm{Ext} }
\newcommand{\PP}{\mathbb{P} }
\newcommand{\ZZ}{\mathbb{Z} }
\newcommand{\bm}{\mathbf{m} }
\newcommand{\bn}{\mathbf{n} }
\newcommand{\bp}{\mathbf{p} }
\newcommand{\cF}{\mathcal{F} }
\newcommand{\cG}{\mathcal{G} }
\newcommand{\cH}{\mathcal{H} }
\def\cHom{\mathcal{H}om }
\newcommand{\cM}{\mathcal{M} }
\newcommand{\cO}{\mathcal{O} }
\newcommand{\cR}{\mathcal{R} }
\newcommand{\cpHom}{\mathcal{P}ar\cHom}
\def\pExt{\mathrm{ParExt} }
\newcommand{\pHom}{\mathrm{ParHom}}
\def\rHom{\mathrm{Hom} }
\title{Ulrich bundles on intersections of two quadrics}
\author{Jiwan Jung, Kyoung-Seog Lee, Han-Bom Moon}
\date{Source: arXiv:2609.09596v1 (CC BY 4.0)}
\begin{document}
\maketitle

\noindent\emph{Source and licence.} Ulrich bundles on intersections of two quadrics. Version arXiv:2609.09596v1 (CC BY 4.0), distributed by arXiv under the Creative Commons Attribution 4.0 International licence, http://creativecommons.org/licenses/by/4.0/, as recorded in the arXiv licence metadata for this exact version and shown on the abstract page https://arxiv.org/abs/2609.09596v1. Full licence notice: \texttt{licenses/arxiv-2609.09596-CC-BY-4.0.txt}.\par\smallskip

\noindent\textit{Editorial scope.} Counted unit: the article's theorem thm:maintheoremoddconditional (source0.tex lines 1226--1231). The article's complete original proof of it is delivered with it (source0.tex lines 1233--1272, one \texttt{proof} environment of 40 lines). No mathematics is written, repaired or re-derived: the statement and the proof are the article's own lines, in the article's own notation, and no retained line is substituted or re-flowed.  Retained uncounted as the source-local context the proof needs: lines 1195--1202.  Nothing was omitted from any retained span, so the package carries no omission record.  No retained line contains a citation, so the wrapper carries no bibliography and no bibliography entry.  No retained line contains a figure environment, an image-inclusion command or drawing code, so no image is delivered and the package declares no asset dependency.  Editorial formatting only, in addition: an AMS \texttt{amsart} preamble that restates the article's own declarations and macros used by the retained lines, namely the environments \texttt{theorem} on separate counters, together with the standard packages those lines need.  The retained environments print in this document as Theorem 1 in order of appearance, whereas the article prints the same items with the numbering of its own full document.  The complete native source of the article as distributed in the arXiv e-print for this exact version is bundled unchanged as \texttt{sources/209834/source0.tex}.
\begin{theorem}\label{thm:initialconst}
\begin{enumerate}
  \item Let \(g \ge 3\) be odd. Let \(\cG_2\) be a parabolic bundle on \((\PP^1, \bp)\) with underlying bundle \(\underline{\cG}_2 = \cO((1-g)/2)^{\oplus 2}\), bit vector \(\bm=(m_i)\) with \(g+1\) nonzero entries, and with a general choice of flags.
  \item Let \(g \ge 3\) be odd. Let \(\cG_3\) be a parabolic bundle on \((\PP^1, \bp)\) with underlying bundle \(\underline{\cG}_3 = \cO((1-g)/2)^{\oplus 3}\), bit vector \(\bn=(n_i)\) with \(3(g+1)/2\) nonzero entries, and a general choice of flags.
  \item Let \(g \ge 2\) be even. Let \(\cG_2^{\mathrm{ev}}\) be a parabolic bundle on \((\PP^1, \bp)\) with underlying bundle \(\underline{\cG}_2^{\mathrm{ev}} = \cO((2-g)/2)\oplus \cO(-g/2)\), bit vector \(\bm=(m_i)\) with \(g+1\) nonzero entries, and a general choice of flags.
\end{enumerate}
Then \(\cG_2\), \(\cG_3\), and \(\cG_2^{\mathrm{ev}}\) are stable and cohomologically orthogonal to \(\cR\).
\end{theorem}

\begin{theorem}\label{thm:maintheoremoddconditional}
\begin{enumerate}
  \item Suppose that \(g \ge 3\) is odd. For each rank \(r \ge 2\), there is an indecomposable vector bundle \(\cF_r\) whose dual is orthogonal to \(\cR\).
  \item Suppose that \(g \ge 2\) is even. For each even \(r \ge 2\), there is an indecomposable vector bundle \(\cF_r\) whose dual is orthogonal to \(\cR\).
\end{enumerate}
\end{theorem}

\begin{proof}
Taking the dual bundle preserves indecomposability. Thus, it is sufficient to construct an indecomposable bundle orthogonal to \(\cR\).

\textbf{Case 1.} \(g\) is odd.

We choose \(\bm\) and \(\bn\) so that \(n_i \ge m_i\) for all \(i\). From the definition of parabolic morphisms, for two parabolic bundles \(\cG := (\underline{\cG}, \{V_i\})\) and \(\cH := (\underline{\cH}, \{W_i\})\), there is a short exact sequence
\begin{align}\label{eq:par_ex_seq}
0\to \cpHom(\cG,\cH)\to\cHom\bigl(\underline{\cG},\underline{\cH}\bigr)\to\bigoplus_{i=1}^{2g+1}\rHom(V_i,\underline{\cH}|_{p_i}/W_i)\otimes k(p_i)\to 0,
\end{align}
which induces a long exact sequence of cohomology groups:
\begin{equation}\label{eqn:par_long_seq}
\begin{split}
  0\to \pHom(\cG,\cH)\to&\rHom\bigl(\underline{\cG},\underline{\cH}\bigr)\to\bigoplus_{i=1}^{2g+1}\rHom(V_i,\underline{\cH}|_{p_i}/W_i)\otimes k(p_i)\\
  &\to \pExt^1(\cG, \cH) \to \Ext^1(\underline{\cG}, \underline{\cH}) \to 0.
\end{split}
\end{equation}
Applying this long exact sequence to stable bundles \(\cG_2\) and \(\cG_3\), we obtain
\[\dim \pExt^1(\cG_2, \cG_2) = g - 2, \quad \dim \pExt^1(\cG_3, \cG_3) = 3g - 5,\]
and
\[\dim \pExt^1(\cG_2, \cG_3) = \dim \pExt^1(\cG_3, \cG_2) = 2g - 4.\]
Since \(g \ge 3\), all extension groups are nontrivial.

For each \(t \in \ZZ_{>0}\), construct \(\cG_{2,t}\) as follows. Set \(\cG_{2,1}:=\cG_2\), and suppose that \(\cG_{2,t-1}\) has been constructed. Then choose \(\cG_{2,t}\) to be a nontrivial extension
\begin{equation}\label{eqn:iteration}
0 \to \cG_2 \to \cG_{2,t} \to \cG_{2,t-1} \to 0.
\end{equation}
There is a nontrivial extension because \(\pExt^1(\cG_{2,t-1},\cG_2)\to\pExt^1(\cG_2,\cG_2)\) is surjective and one may lift a nontrivial extension.

Applying \(\pHom(\cG_2,-)\) to \eqref{eqn:iteration}, we inductively obtain \(\pHom(\cG_2,\cG_{2,t})\cong k\). Since every simple factor of \(\cG_{2,t}\) is isomorphic to \(\cG_2\), this shows that \(\cG_{2,t}\) has a unique simple subobject and hence is indecomposable.

If \(r\) is even, take \(\cG':=\cG_{2,r/2}\). If \(r\) is odd, the same argument applies by replacing the last factor by \(\cG_3\), using \(\pExt^1(\cG_3,\cG_2)\ne0\).

\textbf{Case 2.} \(g\) is even and \(\ge 4\).

In this case, we may apply the same argument. For \(\cG_2^{\mathrm{ev}}\) in Theorem~\ref{thm:initialconst}, using \eqref{eqn:par_long_seq}, we have \(\dim \pExt^1(\cG_2^{\mathrm{ev}}, \cG_2^{\mathrm{ev}}) = g-2\). Therefore, if \(g \ge 4\), the same iterative construction with \(\cG_2^{\mathrm{ev}}\) in place of \(\cG_2\) gives an indecomposable bundle of every even rank \(r\), orthogonal to \(\cR\).

\textbf{Case 3.} \(g = 2\).

In this case, one may choose \(\cG_{2, I}^{\mathrm{ev}} \in \cM_{(\PP^1, \bp)}(2, -1, \bm_I)\) and \(\cG_{2, J}^{\mathrm{ev}} \in \cM_{(\PP^1, \bp)}(2, -1, \bm_J)\), where \(\bm_I = (1, 1, 1, 0, 0)\) and \(\bm_J = (0, 0, 1, 1, 1)\). Both are stable. Using \eqref{eqn:par_long_seq}, we have \(\pExt^1(\cG_{2, I}^{\mathrm{ev}}, \cG_{2, J}^{\mathrm{ev}})\cong\pExt^1(\cG_{2, J}^{\mathrm{ev}}, \cG_{2, I}^{\mathrm{ev}})\cong k\). Starting with either bundle and taking nontrivial extensions successively, alternating \(\cG_{2,I}^{\mathrm{ev}}\) and \(\cG_{2,J}^{\mathrm{ev}}\), the same proof gives an indecomposable bundle of every even rank \(r\), orthogonal to \(\cR\).
\end{proof}

\end{document}
